
\documentclass[11pt,a4paper]{article}

% ==================================================
% 1. PACKAGES
% ==================================================

\usepackage[a4paper, margin=1in]{geometry}
\usepackage{amsmath, amssymb, amsthm, mathtools}
\usepackage{mathrsfs}
\usepackage{bm}
\usepackage{enumitem}
\usepackage{xcolor}
\usepackage[most]{tcolorbox}
\usepackage{hyperref}
\usepackage{cleveref}

\hypersetup{
    colorlinks=true,
    linkcolor=blue,
    citecolor=blue,
    urlcolor=blue
}

\setlength{\parskip}{0.35em}
\setlength{\parindent}{0pt}

% ==================================================
% 2. MATHEMATICAL ENVIRONMENTS
% ==================================================

\theoremstyle{definition}
\newtheorem{definition}{Definition}[section]
\newtheorem{example}[definition]{Example}
\newtheorem{exercise}[definition]{Exercise}
\newtheorem{remark}[definition]{Remark}
\newtheorem{counterexample}[definition]{Counterexample}
\newtheorem{note}[definition]{Note}

\theoremstyle{plain}
\newtheorem{theorem}[definition]{Theorem}
\newtheorem{lemma}[definition]{Lemma}
\newtheorem{proposition}[definition]{Proposition}
\newtheorem{corollary}[definition]{Corollary}

% ==================================================
% 3. PERSONAL NOTE BOX
% ==================================================

\newtcolorbox{mynote}[1][Personal Note]{
    enhanced,
    breakable,
    colback=yellow!5,
    colframe=yellow!50!black,
    title={#1},
    fonttitle=\bfseries,
    boxrule=0.6pt,
    arc=2mm,
    left=2mm,
    right=2mm,
    top=1mm,
    bottom=1mm
}

% ==================================================
% 4. CUSTOM MATHEMATICAL COMMANDS
% ==================================================

\newcommand{\R}{\mathbb{R}}
\newcommand{\N}{\mathbb{N}}
\newcommand{\Q}{\mathbb{Q}}
\newcommand{\Z}{\mathbb{Z}}
\newcommand{\E}{\mathbb{E}}
\newcommand{\A}{\mathcal{A}}
\newcommand{\Prob}{\mathbb{P}}
\newcommand{\eps}{\varepsilon}

\newcommand{\norm}[1]{\left\lVert #1 \right\rVert}
\newcommand{\abs}[1]{\left| #1 \right|}
\newcommand{\inner}[2]{\left\langle #1,#2 \right\rangle}
\newcommand{\set}[1]{\left\{#1\right\}}

% ==================================================
% 5. DOCUMENT INFORMATION
% ==================================================

\title{\textbf{Algebras and $\sigma$-algebras}}

% ==================================================
% 6. DOCUMENT
% ==================================================

\begin{document}

\maketitle
\tableofcontents
\newpage

\section{Motivation}
\begin{mynote}[Why $\sigma$-algebras?]
For further chapters, if we want to construct measures, we will see that we can't define measure for arbitrary set. So we gonna limit those sets to a class, which is $\sigma$-algebras.
\end{mynote}

\section{Definitions, Examples and Properties}
\begin{definition}[Algebra]
An algebra is a collection of $\mathcal{A}$ of subsets of $X$ such that:
\begin{itemize}
    \item \textbf{1.} $\emptyset \in \A$ and $X\in \A$.
    \item \textbf{2.} if $A \in \A$ then $A^c \in \A$.
    \item \textbf{3.} if $A_1, A_2,A_3,..., A_n\in \A$, then the union of $A_i$ and the intersection of $A_i$ are in $\A$.  
\end{itemize}
\end{definition}

\begin{definition}[$\sigma$-algebra]
$\A$ is a $\sigma$-algebra if $\A$ is an algebra and whenever $A_1, A_2,... \in \A$, then the union (may infinite but countable) and the intersection (may infinite but countable) of $A_i$ is in $\A$.

The pair $(X, \A)$ called measurable space. A set $A$ is measurable or $\A$ measurable if $A \in \A$ (note that, $\A$ has to be $\sigma$-algebra, not just  algebra).
\end{definition}

\begin{note}
If we have $2.$ (closed in complement) and closed in union, we can imply closed in intersection.
\end{note}

\begin{mynote}[What to focus?]
The main object in this part is indeed, algebra. Sigma algebra just extends algebra to countable infinite intersection and union.    
\end{mynote}

\begin{mynote}[Our old friend, subspace]
You may see this algebra object as a kind of "subspace" of the power set of $X$ (the set that contains all subset of $X$). To see the "symmetry" (not exactly) in this scenario, we have an $\sigma$-algebra $\A$ of $X$, then:
\begin{itemize}
    \item $\A \subseteq \mathcal{P}(X)$. (subspace is a subset of a vector space)
    \item closed in complement, closed in union (countable infinite or finite). (addition and multiplication are closed in subspace)
    \item contain $X$ (therefore contain $\emptyset$ by closed in complement). (subspace contain $\mathbf{0}$)
\end{itemize}
\end{mynote}

\begin{definition}[Measurable space]
The pair $(X, \A)$ is called a measurable space. A set $A$ is measurable if $A \in \A$.
\end{definition}

\begin{example}[Trivial example]
$X=\R$, $\A$ is the collection of all subsets of $\R$ (power set). Then $\A$ is $\sigma$-algebra. Trivial.
\end{example}

\begin{example}
$X=\R$, $\A = \{A\subset\R: A \text{ is countable or } A^c \text{ is countable} \}$. $\A$ is $\sigma$-algebra.
\end{example}

\begin{proof}

To prove $\A$ is $\sigma$-algebra, we gonna use Definition 2.2.

We have $\emptyset$ is countable, then $\emptyset \in \A$. We have $\R$ which has its complement is $\emptyset$ is countable, then $\R \in \A$. Thus, $\A$ satisfied $(1.)$.

If $A \in \A$, then $A$ is countable or its complement is countable. First case, if $A$ is countable, then $A^c$ has its complement which is $(A^c)^c=A$ is countable, then $A^c\in \A$. Second case, if $A^c$ is countable, then $A^c \in A$. Therefore, $\A$ satisfied $(2.)$.

Consider $A_1, A_2,... \in \A$. First case, if all $A_i$ are countable, then its union will be countable, thus the union will be in $\A$. Second case, if there is a $A_j$ that has $A^c_j$ countable, then we have $(\bigcup_iA_i)^c=\bigcap_iA_i^c\subseteq A_j^c$. But $A_j^c$ is countable, then $(\bigcup_iA_i)^c$ is countable (subset of a countable set is countable). Then, $(\bigcup_iA_i)^c \in \A \Rightarrow \bigcup_iA_i \in \A$. So, $\A$ satisfied (3.).

If you want to prove that the intersection is in $\A$ too, just use closed in complement and close in countable infinite union which was proven above.
\end{proof}


\begin{example}[Closed interval]
Let $X=[0,1]$ and $\A=\{\emptyset, X, [0,\frac{1}{2}], (\frac{1}{2},1]\}$. $\A$ is $\sigma$-algebra. Easy to check.
\end{example}

\begin{example}[Finite set]
$X=\{1,2,3\}$, $\A = \{\emptyset, X, \{1\}, \{2,3\}\}$. $\A$ is $\sigma$-algebra.
\end{example}

\begin{example}
$X=[0,1]$ and $B_1,B_2,..., B_8$, pairwise disjoint, union of all $B_i$ is $X$, let $\A$ be the collection of all finite union of $B_i$. Then $\A$ is $\sigma$-algebra.
\end{example}

\begin{mynote}[$\sigma$-algebra is just coloring segments! (in finite, discrete setting; or see Example 2.8)]
Imagine an array of 6 boxes, it is our $X$. In this example, we gonna use 3 segments. The first segment contains box 1,2; the second segment contains box 3,4,5; and the last segment contains box 6. We have a brush, to color those segment, and remember, we can only color segments, not boxes! Then, $\A$ is \textbf{all of the possible ways} that we can color those segment. For example, we can color segment 1, 3, but leave the segment 2. That's is a way to color those segments, and that way is belong to $\A$! Then, each segment have 2 options, that is no color or color. So the number of elements in $\A$ (finite case) is $2^3=8$. Generally, we have $|\A|=2^k$ with $X$ is finite, and $k$ is the number of segments. Hope this is helpful for you.
\end{mynote}

\begin{lemma}
If $\A_{\alpha}$ is a $\sigma$-algebra for each $\alpha$ in some non-empty index set $I$, then $\bigcap_{\alpha\in I}\A_{\alpha}$ is $\sigma$-algebra.
\end{lemma}

\begin{proof}
This comes directly from Definition 2.1 and 2.2.
\end{proof}

\section{Borel sigma-algebra}

\begin{definition}[Smallest sigma algebra of $\mathcal{C}$]
$\sigma(\mathcal{C}) = \bigcap\{\A_{\alpha}: \A_{\alpha} \text{ is $\sigma$-algebra}, \mathcal{C}\subset \A_{\alpha}\}$
\end{definition}

\begin{note}
$\sigma(\mathcal{C})$ is not empty, we always have the power set in that intersection. By Lemma 2.10, $\sigma(\mathcal{C})$ is $\sigma$-algebra. And $\sigma(\mathcal{C})$ is the smallest sigma-algebra contains $\mathcal{C}$.
\end{note}

\begin{proposition}
Here are some properties of $\sigma(\cdot)$:
\begin{itemize}
    \item \textbf{Monotonic:} $B\subset C$ then $\sigma(B) \subset \sigma(C)$.
    \item \textbf{Idempotent:} $\sigma(\sigma(C)) = \sigma(C)$.
\end{itemize}
\end{proposition}

\begin{definition}[Borel $\sigma$-algebra]
Given $(X, d)$ is a metric space, let $\mathcal{G}$ is the set of all open set in $X$. Then $\mathcal{B} = \sigma(\mathcal{G})$ is Borel $\sigma$-algebra on $X$, elements of $\mathcal{B}$ is called Borel set, or Borel measurable.
\end{definition}

\begin{mynote}[Smallest sigma vs. span, why measure theory?]
The smallest sigma algebra is "somehow" the same with span of a set of vector in linear algebra. The smallest sigma embeds an initial collection of subset to a "subspace" that preserves Boolean operator (union, intersection, complement), and that "subspace" is the smallest (in size, maybe).

There is a note in the book, which is: In case $X$ is real line, $\mathcal{B} \neq \mathcal{P}(\R)$, it means that, not every subset of $\R$ is Borel measurable. That's why we need measure theory.
\end{mynote}

\begin{proposition}[Ways to generate $\mathcal{B}$, proof omitted]
If $X = \R$, then $\mathcal{B}$ is generated by:
\begin{itemize}
    \item $C_1 =\{(a,b), a,b\in \R\}$.
    \item $C_2 =\{[a,b], a,b\in \R\}$.
    \item $C_3 =\{(a,b], a,b\in \R\}$.
    \item $C_4 =\{(a,\infty), a\in \R\}$.
\end{itemize}
\end{proposition}

\begin{note}
You can see that in Proposition 3.5, you still can use closed intervals for creating Borel sigma-algebra (because Borel sigma-algebra still contains closed sets!).
\end{note}

\section{Exercises}

All exercises in this note are from "Real analysis for Graduate Students" by Richard F. Bass.

\begin{exercise}
\textit{\textbf{Find an example of a set $X$ and two sigma-algebras $A_1$ and $A_2$, each consisting of subsets of $X$, such that $A_1 \cup A_2$ is not a sigma-algebra.}}

$X = \{1,2,3\}, A_1=\{\emptyset, X, \{1\}, \{2,3\}\}, A_2=\{\emptyset, X, \{1,2\}, \{3\}\}$.
Union of $A_1, A_2$ violates closed in intersection, $\{1,2\} \cap \{2,3\} = \{2\}\not\in (A_1\cup \A_2)$.
\end{exercise}

\begin{exercise}
\textbf{\textit{See Exercise 2.2 p.15}}

Let $X = \N$, and $\A_i$ is the $\sigma$-algebra generated by the partition $1-2-3-...-i-i+1,i+2,...$. Thus, $\A_i$ will consist of finite unions of finite unions from $\{1\}, \{2\}, \{3\},...,\{i\}$ and $\{i+1,i+2,i+3,...\}$.

For example, let generate $\A_2$, which is from the partition $1-2-3,4,5...$. We have $\A_2=\{\emptyset, \N, \{1\}, \{2\}, \{1,2\}, R_2, R_2\cup \{1\}, R_2\cup \{2\}\}$, with $R_i = \{i+1,i+2,i+3,...\}$ and $i \in \N$.

We want to proof that $\A_i \subset \A_{i+1}$. So let $x \in \A_i$. If $x$ is in the finite unions of $\{1\}, \{2\}, \{3\},...,\{i\}$ which is a subset of finite unions of $\{1\}, \{2\}, \{3\},...,\{i\}, \{i+1\}$, then $x \in \A_{i+1}$. If $x$ is in the finite unions from $\{1\}, \{2\}, \{3\},...,\{i\}$ and $\{i+1,i+2,i+3,...\}$ which is a subset of the finite unions from $\{1\}, \{2\}, \{3\},...,\{i+1\}$ and $\{i+2,i+3,i+4,...\}$ (because from the partition of $\A_{i+1}$, we can recreate $R_i$ by just union $R_{i+1}$ with $\{i+1\}$, after that we obtain $\A_{i}$). Thus $\A_{i}\subset \A_{i+1}$.

Now we have already set up the initial condition of the exercise. For $\bigcup_i^{\infty}\A_i$ not to be a $\sigma$-algebra, we need to violate condition \textbf{3.} in Definition 2.1, which is closed in union (may countably infinite).

Let $E= \{2,4,6,8,...\}$, which is the union of $\{2\}, \{4\},\{6\},...$ and $\{2k\} \in \bigcup_i^{\infty}\A_i$. Now we need to show that, $E \not\in \bigcup_i^{\infty}\A_i$ to violate \textbf{3.}. Assume $E \in \bigcup_i^{\infty}\A_i$. For a $I \in \bigcup_i^{\infty}\A_i$, it must be in some $\A_j$. Thus, any $I \in \bigcup_i^{\infty}\A_i$ must be in one of those 2 form: contain every elements of $R_j$ or contain no elements of $R_j$. Because $E \in \bigcup_i^{\infty}\A_i$, WLOG (we can extend to multiple $j$), assume there is a only $j$ such that $E \in \A_j$. If $E$ contain every elements of $R_j$, it is impossible, because $R_j$ contains odd numbers. If $E$ contain no elements of $R_j$, then $E$ will miss all even number after $j$. Therefore, $E \not\in \A_j \Rightarrow E \not\in \bigcup_i^{\infty}\A_i$, which is a contradiction. Then, $ E \not\in \bigcup_i^{\infty}\A_i$, violate condition \textbf{3.}. Thus, $\bigcup_i^{\infty}\A_i$ is not a $\sigma$-algebra in this case.
\end{exercise}

\begin{exercise}
\textbf{\textit{See Exercise 2.5 p.16}}

Before jumping into this exercise, you have to know that, $f^{-1}$ is not the inverse function of $f$ (because $f$ is not bijective!). In this case, $f^{-1}$ is the "inverse image function", which map $\mathcal{P}(Y) \rightarrow \mathcal{P}(X)$ and $f^{-1}(y) = \{x \in \mathcal{P}(X): f(x) = y\}$.

For example, we have $X = \{1,2,3,4,5\}$, $f: X \rightarrow Y, Y=\{a,b,c,d,e,f\}, f(1) = a, f(2) = b, f(3) = c, f(4) = d, f(5) = d$. And with $f^{-1}: \mathcal{P}(Y) \rightarrow \mathcal{P}(X)$, I can just give you some examples of its result: $f^{-1}(\emptyset) = \emptyset, f^{-1}(\{a\}) = \{1\}, f^{-1}(\{b\}) = \{2\}, f^{-1}(\{c\}) = \{3\}, f^{-1}(\{a,c\}) = \{1, 3\}, f^{-1}(\{a,d,f\}) = \{1,4,5\}, f^{-1}(\{e,f\}) = \{\emptyset\}, f^{-1}(\{Y\}) = \{X\}$.

After getting some senses what a reverse image function doing, we can go to the proof, of course by using 3 condition in Definition 2.1, 2.2.

\begin{proof}
(Using definition):

\textbf{Contains empty set and $X$.} As we discussed above, we have $f^{-1}(\emptyset) = \emptyset$ and $f^{-1}(Y) = X$. Therefore, $\emptyset, X \in \mathcal{B} = \{f^{-1}(A): A \in \A \}$.

\textbf{Closed in complement.} Let $B_0 \in \mathcal{B}$, then there is a $A \in \A$ s.t. $f^{-1}(A) = B_0$. By the property of $f^{-1}$, we have, $f^{-1}(A^c) = X\setminus B_0$. Then, $X\setminus B_0 = (B_0)^c\in \mathcal{B}$.

\textbf{Closed in countable infinite union.} Let $B_1, B_2,... \in \mathcal{B}$. We have $\bigcup_i^{\infty}B_i = \bigcup_i^{\infty}f^{-1}(A_i) = f^{-1}(\bigcup_i^{\infty}A_i) \in \mathcal{B}$.
\end{proof}
\end{exercise}

\end{document}
